% \begin{table}[t]
% \centering
% \footnotesize
% \setlength{\tabcolsep}{6pt}
% \begin{threeparttable}
% \caption{%
% Architecture ablations on DINOv2 ViT-B/14.
% Each row modifies one component from the full \model{} configuration.
% % \MT{PODS is 58.3 I think .. there may be typos in this table and then corresponding text}
% % \MG{No 58.3 is with MLLM descriptions, 53.4 is for human generated ones}
% }
% \label{tab:arch_ablation}
% \begin{tabular}{@{}l
%                 S[table-format=2.1]
%                 S[table-format=2.1]
%                 S[table-format=2.1]@{}}
% \toprule
% {Variant}
%   & {FG-CLS $\uparrow$}
%   & {Steer.\ $\uparrow$}
%   & {PODS $\uparrow$} \\
% \midrule
% \color{gray} DINOv2 (frozen)
%   & \color{gray} 89.0
%   & \color{gray} 43.5
%   & \color{gray} 28.6 \\
% \rowcolor{\methodcolor}
% \textbf{\model{} (ours)}
%    & \textbf{76.3}
%    & \textbf{95.9}
%    & \textbf{53.4} \\
% \addlinespace
% \rowcolor{gray!20}
% \multicolumn{4}{l}{Architecture ablations} \\
% Late fusion
%   & 86.5
%   & 93.4
%   & 45.3 \\
% w/o tanh gating
%   & 36.2
%   & 28.4
%   & 10.8 \\
% Linear text projector
%   & 71.5
%   & 94.2
%   & 51.5 \\
% % w/o FFN
% %   & 76.3
% %   & 95.9
% %   & 53.4 \\
% \addlinespace
% % \multicolumn{4}{l}{\YA{"Location of cross attn layer"; ADD THESE LINES/subheadings through table}} \\
% \rowcolor{gray!20}
% % \multicolumn{4}{l}{Location of cross-attention layers} \\
% % Early layers (1, 3)
% %   & {XX.X} & {XX.X} & {XX.X} \\
% % Early+mid (1, 3, 5, 7)
% %   & {XX.X} & {XX.X} & {XX.X} \\
% % Mid+late (5, 7, 9 11)
% %   & {XX.X} & {XX.X} & {XX.X} \\
% % Late layers (9, 11)
% %   & {XX.X} & {XX.X} & {XX.X} \\
% % \addlinespace
% % Gaussian point GT
% %   & {XX.X} & {XX.X} & {XX.X} \\
% \bottomrule
% \end{tabular}
% \begin{tablenotes}[flushleft]\footnotesize
% \item Late fusion appends gated cross-attention after the final ViT layer instead of interleaving within layers.
% PODS is evaluated with descriptive text prompts; the frozen baseline uses no text. Both tanh gating on the cross attention blocks and MLP projector RoBERTa are detrimental to model performance.
% % Layer-placement rows vary which ViT blocks receive cross-attention (default: blocks 2, 5, 8, 11).
% % Gaussian point GT replaces the segmentation-mask target with a Gaussian heatmap centered on the bounding-box midpoint ($\sigma{=}1.1$).
% \end{tablenotes}
% \end{threeparttable}
% \end{table}

\begin{wraptable}{r}{0.50\textwidth}
% \begin{table}[t]
\centering
\footnotesize
\tabcolsep=0.5mm
\begin{threeparttable}
\vspace{-9mm}
\caption{%
Architecture ablations on DINOv2 ViT-B/14.
\model{} combines three design choices; each ablation row toggles exactly one (marked \xmark).
}
\vspace{-5mm}
\label{tab:arch_ablation}
\begin{tabular}{@{}ccc
                S[table-format=2.1]
                S[table-format=2.1]
                S[table-format=2.1]@{}}
\toprule
{Early} & {Tanh} & {MLP} & {FG-CLS} & {CORE} & {PODS} \\
{Fusion} & {Gate}  & {Proj.} & $\uparrow$ & $\uparrow$ & $\uparrow$ \\
\midrule
\color{gray} {--}
  & \color{gray} {--}
  & \color{gray} {--}
  & \color{gray} 89.0
  & \color{gray} 43.5
  & \color{gray} 28.6 \\
\rowcolor{\methodcolor}
\cmark & \cmark & \cmark
  & 87.7
  & \textbf{96.0}
  & \textbf{58.1} \\
\addlinespace
\xmark & \cmark & \cmark
  & \textbf{91.8}
  & 93.3
  & 36.6 \\
\cmark & \xmark & \cmark
  & 83.5
  & 94.6
  & 47.1 \\
\cmark & \cmark & \xmark
  & 86.7
  & 95.2
  & 56.4 \\
\bottomrule
\end{tabular}
% \begin{tablenotes}[flushleft]\footnotesize
Top gray row: frozen DINOv2 baseline (no adapter \JR{in Sec. 3, we call the MLP adapter; might be better to use other term here.}).
\xmark\ in Early Fusion means late fusion. %(cross-attention appended after the final ViT layer).
\xmark\ in Tanh Gate means ungated cross-attention.
\xmark\ in MLP Proj.\ means single linear layer instead of two. % replacing the two-layer MLP.
PODS uses descriptive text prompts.
Frozen DINOv2 baseline with no adapter (top row, gray) cannot use text. \JR{merge with first sentence}
% \end{tablenotes}
\end{threeparttable}
\vspace{-6mm}
% \end{table}
\end{wraptable}

